\begin{Problem}
	\textbf{Exercise 4.1: Convolution product \hfill (6P)}
	
	Let us consider the probability density of the gaps between Poisson-distributed events
	(the length $x$ of silence between the clicks of a Geiger counter), which is defined by
	\[
	p(x)=w e^{-wx}
	\qquad (x\ge 0,\; w>0).
	\]
	
	\begin{enumerate}[label=(\alph*)]
		\item Compute the corresponding characteristic function (see lecture notes). \hfill (1P)
		
		\item Compute the two-fold convolution product $(p*p)(x)$ and the three-fold convolution
		product $(p*p*p)(x)$. \hfill (2P)
		
		\item Guess a general formula for the $N$-fold convolution product
		$p^{*N}(x)$ and prove it by recursion. \hfill (2P)
		
		\item Fit a normal distribution with mean $\mu$ and width $\sigma$ to the probability
		density computed in (c) for large $N$.
		
		\textit{Hint:} Recall the behavior of cumulants. \hfill (1P)
	\end{enumerate}
\end{Problem}

\begin{Problem}
	\textbf{Exercise 4.2: Harmonic oscillator \hfill (6P)}
	
	Consider a harmonic oscillator with eigenstates
	\[
	|n\rangle = |0\rangle, |1\rangle, |2\rangle, \ldots
	\]
	and eigenenergies
	\[
	E_n = \hbar \omega \left(n+\frac12\right).
	\]
	
	Let us assume that the eigenstates correspond to classical configurations and that
	the system evolves stochastically according to a continuous-time Markov process as follows:
	\[
	n \to n+1 \qquad \text{with rate } \omega_\uparrow
	\]
	\[
	n \to n-1 \qquad \text{with rate } \omega_\downarrow
	\]
	
	\begin{enumerate}[label=(\alph*)]
		\item Calculate the Liouville operator $\mathcal{L}$. \hfill (1P)
		
		\item Determine the left and the right eigenvector to the eigenvalue zero. \hfill (2P)
		
		\item Normalize the right eigenvector obtained in (b) so that it can be interpreted
		as a probability distribution. Under which condition is the normalization possible?
		\hfill (2P)
		
		\item In a canonical ensemble the probability of finding the system at energy $E_n$
		is proportional to
		\[
		\exp\!\left(-\frac{E_n}{k_B T}\right),
		\]
		where $T$ is the temperature and $k_B$ is the Boltzmann constant.
		Express the temperature $T$ in terms of the rates
		$\omega_\uparrow$ and $\omega_\downarrow$. \hfill (1P)
	\end{enumerate}
\end{Problem}